\slajdid{10-s015}
\begin{frame}[label=10-s015]
\gusttekst\setlength{\abovedisplayskip}{3pt}\setlength{\belowdisplayskip}{3pt}\setlength{\abovedisplayshortskip}{2pt}\setlength{\belowdisplayshortskip}{2pt}\[V_O=V_{CC}-\beta_F\frac{R_C}{R_B}(V_I-0.6V)\]
Namerno je napisano u ovim oblicima pošto ovde mogu da nastanu nedoumice da li smo dobro uradili. Odnosno da neki rezultati budu zbunjujući.
\par\medskip
Ako bi hteli da nastavimo da crtamo naš grafik, očigledno je da moramo krenuti od tačke gde smo stali, odnosno za $V_I=V_{\gamma T}$. Međutim ako ovu tačku zamenimo u jednačini koju smo prethodno našli
\[V_O=V_{CC}-\beta_F\frac{R_C}{R_B}(V_{\gamma T}-0.6V)=V_{CC}-\beta_F\frac{R_C}{R_B}(0.5V-0.6V)>V_{CC}\quad???\]
Ovakav rezultat možda jeste nelogičan, ali je očekivan zbog naše diskretne prirode modela tranzistora i „naglog skoka“ sa $V_{BE1}=V_{\gamma T}=0.5V\ na\ V_{BE1}=V_{BE}=0.6V$. Ovakve „nelogičnosti“ će nam se i dalje pojavljivati, međutim greške koje pravimo sa našim diskretnim modelima su male i nastavićemo da radimo sa njima. Znači smatraćemo da nema prekida u tački $V_I=V_{\gamma T}$ da se jednostavno od te tačke nastavlja karakteristika prenosa po jednačini
\[V_O=V_{CC}-\beta_F\frac{R_C}{R_B}(V_I-V_{BE})\]
\end{frame}
